% exam_question_ot_kl.tex

\documentclass[12pt]{article}
\usepackage{amsmath,amssymb}
\usepackage[margin=1.2in]{geometry}
\usepackage{setspace}
\usepackage{parskip}
\onehalfspacing

\begin{document}

\begin{center}
  {\large\bfseries Exam Question: Couplings, KL Divergence, and Optimal Transport}
\end{center}

\bigskip

Let $X \subset \mathbb{R}$ be a compact interval and let $f$ and $g$ be two
distinct probability density functions defined on $X$, both positive everywhere
on $X$ and integrating to one with respect to Lebesgue measure~$\lambda$.
\begin{enumerate}
  \item[\textbf{(a)}] \textbf{Coupling.}  Define a \emph{coupling} of $f$ and $g$.
    In particular, specify what joint density $\pi(x,y)$ on $X \times X$ must
    satisfy in order to qualify as a coupling, and give one explicit example.

  \item[\textbf{(b)}] \textbf{KL divergence.}  Write down the
    Kullback-Leibler divergence $D_{\mathrm{KL}}(f \| g)$.
    State (with brief justification) whether it is symmetric in $f$ and $g$,
    and whether it constitutes a metric on the space of density functions.

  \item[\textbf{(c)}] \textbf{Optimal transport problem.}  Let $c: X \times X
    \to [0,\infty)$ be a continuous cost function.  Recalling the definition
    of a coupling from part~(a), write down the Kantorovich optimal transport
    problem that seeks a coupling of $f$ and~$g$ that minimizes expected
    transport cost.

  \item[\textbf{(d)}] \textbf{Wasserstein distance.}  For $p \geq 1$, define
    the Wasserstein-$p$ distance $W_p(f,g)$ as the outcome of a  special case of the
    Kantorovich problem in~(c), specifying the choice of cost function.
    State one property that $W_p$ possesses that $D_{\mathrm{KL}}(f\|g)$ lacks.
\end{enumerate}

\end{document}
